% GENERATED FILE. Edit canonical lessons, then run npm run generate:books.
% canonical-source-hash: fnv1a64:cfd085d5ca6b62a2
% canonical-lessons: TE-C540-sahityam, TE-C540-kathanika, TE-C540-racana, TE-C540-racincu, TE-C540-anuvadam

\chapter{Literature}
\label{ch:te-literature}
\hlchaptermodality{540}

\begin{chapteropening}
\textbf{By the end of this chapter:} \emph{I can talk about literature, short stories, writing and translation.}

Complete the last lesson of chapter 540: I can talk about literature, short stories, writing and translation.
\end{chapteropening}

\section[\texorpdfstring{\te{సాహిత్యం} (sāhityaṁ)}{సాహిత్యం (sāhityaṁ)}]{\te{సాహిత్యం} (sāhityaṁ) — literary works; literature}

\label{lesson:TE-C540-sahityam}



\begin{quote}
Before the new one: say the Telugu for a lyric, then the Telugu for folk tradition.
\end{quote}

\subsection*{You'll want to know: \te{సాహిత్యం}}
\textbf{\te{సాహిత్యం}} — \emph{sāhityaṁ} — \textquotedblleft{}literary works; literature\textquotedblright{}.

\begin{quote}
\te{తెలుగు} \te{సాహిత్యం} \te{చాలా} \te{పాతది}.
\end{quote}

\begin{quote}
\emph{telugu sāhityaṁ cālā pātadi.}
\end{quote}

\begin{quote}
Telugu literature is very old.
\end{quote}

The everyday word; \te{వాఙ్మయం} is more learned.

\begin{cousinweb}
\textbf{\te{సాహిత్యం}} (\emph{sāhityaṁ}) is borrowed from Sanskrit \emph{sāhitya}, \textquotedblleft{}association, union; literary composition\textquotedblright{} (Monier-Williams).
\end{cousinweb}

\begin{grammarlens}[title={What you've built}]
One more word for describing what you lived and felt.
\end{grammarlens}

\subsection*{Your turn}
\begin{itemize}
\raggedright
  \item \emph{Say it:} \emph{sāhityaṁ}
  \item \emph{Say it:} \emph{sāhityaṁ}, once more
  \item \emph{Say it:} \emph{telugu sāhityaṁ cālā pātadi.} — Telugu literature is very old.
  \item \emph{Recall:} say the Telugu for a lyric, then the Telugu for folk tradition, then say \emph{sāhityaṁ} again
\end{itemize}

\begin{tcolorbox}[breakable,colback=teal!4,colframe=teal!35!black,title={Before you move on}]
What does \te{సాహిత్యం} mean? (\textquotedblleft{}Literary works; literature\textquotedblright{}.) Say it once more.
\end{tcolorbox}
\section[\texorpdfstring{\te{కథానిక} (kathānika)}{కథానిక (kathānika)}]{\te{కథానిక} (kathānika) — a short story}

\label{lesson:TE-C540-kathanika}



\begin{quote}
Before the new one: say the Telugu for folk tradition, then the Telugu for literary works.
\end{quote}

\subsection*{You'll want to know: \te{కథానిక}}
\textbf{\te{కథానిక}} — \emph{kathānika} — \textquotedblleft{}a short story\textquotedblright{}.

\begin{quote}
\te{ఆ} \te{పత్రికలో} \te{ఆమె} \te{కథానిక} \te{వచ్చింది}.
\end{quote}

\begin{quote}
\emph{ā patrikalō āme kathānika vaccindi.}
\end{quote}

\begin{quote}
Her short story appeared in that magazine.
\end{quote}

Shorter than a \te{కథ}: the modern short-story form.

\begin{cousinweb}
\textbf{\te{కథానిక}} (\emph{kathānika}) is built on Sanskrit \emph{kathā}, \textquotedblleft{}conversation; story\textquotedblright{} (CDIAL 2704), with a diminutive ending.
\end{cousinweb}

\begin{grammarlens}[title={What you've built}]
One more word for describing what you lived and felt.
\end{grammarlens}

\subsection*{Your turn}
\begin{itemize}
\raggedright
  \item \emph{Say it:} \emph{kathānika}
  \item \emph{Say it:} \emph{kathānika}, once more
  \item \emph{Say it:} \emph{ā patrikalō āme kathānika vaccindi.} — Her short story appeared in that magazine.
  \item \emph{Recall:} say the Telugu for folk tradition, then the Telugu for literary works, then say \emph{kathānika} again
\end{itemize}

\begin{tcolorbox}[breakable,colback=teal!4,colframe=teal!35!black,title={Before you move on}]
What does \te{కథానిక} mean? (\textquotedblleft{}A short story\textquotedblright{}.) Say it once more.
\end{tcolorbox}
\section[\texorpdfstring{\te{రచన} (racana)}{రచన (racana)}]{\te{రచన} (racana) — a piece of writing; a work}

\label{lesson:TE-C540-racana}



\begin{quote}
Before the new one: say the Telugu for literary works, then the Telugu for a short story.
\end{quote}

\subsection*{You'll want to know: \te{రచన}}
\textbf{\te{రచన}} — \emph{racana} — \textquotedblleft{}a piece of writing; a work\textquotedblright{}.

\begin{quote}
\te{ఆ} \te{రచయిత} \te{కొత్త} \te{రచన} \te{బాగా} \te{అమ్ముడుపోయింది}.
\end{quote}

\begin{quote}
\emph{ā racayita kotta racana bāgā ammuḍupōyindi.}
\end{quote}

\begin{quote}
That writer's new work sold very well.
\end{quote}

Plural: \te{రచనలు}, \emph{racanalu}.

Plural \te{రచనలు}, a writer's works.

\begin{cousinweb}
\textbf{\te{రచన}} (\emph{racana}) is borrowed from Sanskrit \emph{racanā}, \textquotedblleft{}arrangement; a literary production, a composition\textquotedblright{} (Monier-Williams).
\end{cousinweb}

\begin{grammarlens}[title={What you've built}]
One more word for describing what you lived and felt.
\end{grammarlens}

\subsection*{Your turn}
\begin{itemize}
\raggedright
  \item \emph{Say it:} \emph{racana}
  \item \emph{Say it:} \emph{racana}, once more
  \item \emph{Say it:} \emph{ā racayita kotta racana bāgā ammuḍupōyindi.} — That writer's new work sold very well.
  \item \emph{Recall:} say the Telugu for literary works, then the Telugu for a short story, then say \emph{racana} again
\end{itemize}

\begin{tcolorbox}[breakable,colback=teal!4,colframe=teal!35!black,title={Before you move on}]
What does \te{రచన} mean? (\textquotedblleft{}A piece of writing; a work\textquotedblright{}.) Say it once more.
\end{tcolorbox}
\section[\texorpdfstring{\te{రచించు} (raciñcu)}{రచించు (raciñcu)}]{\te{రచించు} (raciñcu) — to compose; to author}

\label{lesson:TE-C540-racincu}



\begin{quote}
Before the new one: say the Telugu for a short story, then the Telugu for a piece of writing.
\end{quote}

\subsection*{You'll want to know: \te{రచించు}}
\textbf{\te{రచించు}} — \emph{raciñcu} — \textquotedblleft{}to compose; to author\textquotedblright{}.

\begin{quote}
\te{ఆ} \te{కవి} \te{వంద} \te{పాటలు} \te{రచించారు}.
\end{quote}

\begin{quote}
\emph{ā kavi vanda pāṭalu raciñcāru.}
\end{quote}

\begin{quote}
That poet composed a hundred songs.
\end{quote}

Past (I …): \te{రచించాను}, \emph{raciñcānu}.

Formal, for books, songs and plays.

\begin{cousinweb}
\textbf{\te{రచించు}} (\emph{raciñcu}) is built on Sanskrit \emph{racana}, \textquotedblleft{}the act of making, arranging, composing\textquotedblright{} (Monier-Williams), with the Telugu verb ending -\te{ఇంచు}.
\end{cousinweb}

\begin{grammarlens}[title={What you've built}]
One more word for describing what you lived and felt.
\end{grammarlens}

\subsection*{Your turn}
\begin{itemize}
\raggedright
  \item \emph{Say it:} \emph{raciñcu}
  \item \emph{Say it:} \emph{raciñcu}, once more
  \item \emph{Say it:} \emph{ā kavi vanda pāṭalu raciñcāru.} — That poet composed a hundred songs.
  \item \emph{Recall:} say the Telugu for a short story, then the Telugu for a piece of writing, then say \emph{raciñcu} again
\end{itemize}

\begin{tcolorbox}[breakable,colback=teal!4,colframe=teal!35!black,title={Before you move on}]
What does \te{రచించు} mean? (\textquotedblleft{}To compose; to author\textquotedblright{}.) Say it once more.
\end{tcolorbox}
\section[\texorpdfstring{\te{అనువాదం} (anuvādaṁ)}{అనువాదం (anuvādaṁ)}]{\te{అనువాదం} (anuvādaṁ) — a translation}

\label{lesson:TE-C540-anuvadam}



\begin{quote}
Before the new one: say the Telugu for a piece of writing, then the Telugu for to compose.
\end{quote}

\subsection*{You'll want to know: \te{అనువాదం}}
\textbf{\te{అనువాదం}} — \emph{anuvādaṁ} — \textquotedblleft{}a translation\textquotedblright{}.

\begin{quote}
\te{ఆ} \te{నవలకు} \te{తెలుగు} \te{అనువాదం} \te{వచ్చింది}.
\end{quote}

\begin{quote}
\emph{ā navalaku telugu anuvādaṁ vaccindi.}
\end{quote}

\begin{quote}
A Telugu translation of that novel has come out.
\end{quote}

\te{అనువాదం} \te{చేయు}, to translate.

\begin{cousinweb}
\textbf{\te{అనువాదం}} (\emph{anuvādaṁ}) is borrowed from Sanskrit \emph{anuvāda}, \textquotedblleft{}saying again, repetition; translation\textquotedblright{} (Monier-Williams).
\end{cousinweb}

\begin{grammarlens}[title={What you've built}]
One more word for describing what you lived and felt.
\end{grammarlens}

\subsection*{Your turn}
\begin{itemize}
\raggedright
  \item \emph{Say it:} \emph{anuvādaṁ}
  \item \emph{Say it:} \emph{anuvādaṁ}, once more
  \item \emph{Say it:} \emph{ā navalaku telugu anuvādaṁ vaccindi.} — A Telugu translation of that novel has come out.
  \item \emph{Recall:} say the Telugu for a piece of writing, then the Telugu for to compose, then say \emph{anuvādaṁ} again
\end{itemize}

\begin{tcolorbox}[breakable,colback=teal!4,colframe=teal!35!black,title={Before you move on}]
What does \te{అనువాదం} mean? (\textquotedblleft{}A translation\textquotedblright{}.) Say it once more.
\end{tcolorbox}
